% GENERATED FILE. Edit canonical lessons, then run npm run generate:books.
% canonical-source-hash: fnv1a64:5fd7ff2fcc00d0aa
% canonical-lessons: MW-C124-kagsi, MW-C124-saban, MW-C124-tauliyo, MW-C124-sui, MW-C124-doro

\chapter{Comb, Soap, Towel, Needle, Thread}
\label{ch:mw-comb-soap-towel-needle-thread}
\hlchaptermodality{124}

\begin{chapteropening}
\textbf{By the end of this chapter:} \emph{I can say a comb, soap, a towel, a needle and thread.}

Complete the last lesson of chapter 124: I can say a comb, soap, a towel, a needle and thread.
\end{chapteropening}

\section[\texorpdfstring{\mw{कांगसी} (kā̃gsī)}{कांगसी (kā̃gsī)}]{\mw{कांगसी} (kā̃gsī) — a comb}

\label{lesson:MW-C124-kagsi}



\begin{quote}
Before the new one: say the Marwadi for a curtain, then the Marwadi for a mirror.
\end{quote}

\subsection*{You'll want to know: \mw{कांगसी}}
\textbf{\mw{कांगसी}} — \emph{kā̃gsī} — \textquotedblleft{}a comb\textquotedblright{}.

\begin{cousinweb}
\textbf{\mw{कांगसी}} (\emph{kā̃gsī}) is a tadbhava, a word inherited from Sanskrit and reshaped by sound change. It comes from \emph{kaṅkata}, \textquotedblleft{}comb\textquotedblright{}. CDIAL gives no Marwari form in this entry (CDIAL 2598); the match is through Gujarati \emph{kā̃gsī}. CDIAL notes that the Gujarati form swapped its sounds, and calls the further origin of the Sanskrit word doubtful.
\end{cousinweb}

\begin{grammarlens}[title={What you've built}]
One more word for this chapter.
\end{grammarlens}

\subsection*{Your turn}
\begin{itemize}
\raggedright
  \item \emph{Say it:} \emph{kā̃gsī}
  \item \emph{Say it:} \emph{kā̃gsī}, once more
  \item \emph{Say it:} say \emph{kā̃gsī}
  \item \emph{Recall:} say the Marwadi for a curtain, then the Marwadi for a mirror, then say \emph{kā̃gsī} again
\end{itemize}

\begin{tcolorbox}[breakable,colback=teal!4,colframe=teal!35!black,title={Before you move on}]
What does \mw{कांगसी} mean? (\textquotedblleft{}A comb\textquotedblright{}.) Say it once more.
\end{tcolorbox}
\section[\texorpdfstring{\mw{साबण} (sābaṇ)}{साबण (sābaṇ)}]{\mw{साबण} (sābaṇ) — soap}

\label{lesson:MW-C124-saban}



\begin{quote}
Before the new one: say the Marwadi for a mirror, then the Marwadi for a comb.
\end{quote}

\subsection*{You'll want to know: \mw{साबण}}
\textbf{\mw{साबण}} — \emph{sābaṇ} — \textquotedblleft{}soap\textquotedblright{}.

\begin{cousinweb}
\textbf{\mw{साबण}} (\emph{sābaṇ}) is a loan, usually traced to Arabic \emph{ṣābūn}, \textquotedblleft{}soap\textquotedblright{}, which reached Marwadi through Persian; its Arabic root is \emph{ṣ-b-n} (Lane). Portuguese \emph{sabão} is another proposed source.
\end{cousinweb}

\begin{grammarlens}[title={What you've built}]
One more word for this chapter.
\end{grammarlens}

\subsection*{Your turn}
\begin{itemize}
\raggedright
  \item \emph{Say it:} \emph{sābaṇ}
  \item \emph{Say it:} \emph{sābaṇ}, once more
  \item \emph{Say it:} say \emph{sābaṇ}
  \item \emph{Recall:} say the Marwadi for a mirror, then the Marwadi for a comb, then say \emph{sābaṇ} again
\end{itemize}

\begin{tcolorbox}[breakable,colback=teal!4,colframe=teal!35!black,title={Before you move on}]
What does \mw{साबण} mean? (\textquotedblleft{}Soap\textquotedblright{}.) Say it once more.
\end{tcolorbox}
\section[\texorpdfstring{\mw{तौलियो} (tauliyo)}{तौलियो (tauliyo)}]{\mw{तौलियो} (tauliyo) — a towel}

\label{lesson:MW-C124-tauliyo}



\begin{quote}
Before the new one: say the Marwadi for a comb, then the Marwadi for soap.
\end{quote}

\subsection*{You'll want to know: \mw{तौलियो}}
\textbf{\mw{तौलियो}} — \emph{tauliyo} — \textquotedblleft{}a towel\textquotedblright{}.

\begin{cousinweb}
\textbf{\mw{तौलियो}} (\emph{tauliyo}) is a loan from Portuguese \emph{toalha}, \textquotedblleft{}towel\textquotedblright{}.
\end{cousinweb}

\begin{grammarlens}[title={What you've built}]
One more word for this chapter.
\end{grammarlens}

\subsection*{Your turn}
\begin{itemize}
\raggedright
  \item \emph{Say it:} \emph{tauliyo}
  \item \emph{Say it:} \emph{tauliyo}, once more
  \item \emph{Say it:} say \emph{tauliyo}
  \item \emph{Recall:} say the Marwadi for a comb, then the Marwadi for soap, then say \emph{tauliyo} again
\end{itemize}

\begin{tcolorbox}[breakable,colback=teal!4,colframe=teal!35!black,title={Before you move on}]
What does \mw{तौलियो} mean? (\textquotedblleft{}A towel\textquotedblright{}.) Say it once more.
\end{tcolorbox}
\section[\texorpdfstring{\mw{सुई} (suī)}{सुई (suī)}]{\mw{सुई} (suī) — a needle}

\label{lesson:MW-C124-sui}



\begin{quote}
Before the new one: say the Marwadi for soap, then the Marwadi for a towel.
\end{quote}

\subsection*{You'll want to know: \mw{सुई}}
\textbf{\mw{सुई}} — \emph{suī} — \textquotedblleft{}a needle\textquotedblright{}.

\begin{cousinweb}
\textbf{\mw{सुई}} (\emph{suī}) is a tadbhava, a word inherited from Sanskrit and reshaped by sound change. It comes from \emph{sūcī}, \textquotedblleft{}needle\textquotedblright{}. CDIAL gives no Marwari form in this entry (CDIAL 13551); the match is through Hindi \emph{sūī}.
\end{cousinweb}

\begin{grammarlens}[title={What you've built}]
One more word for this chapter.
\end{grammarlens}

\subsection*{Your turn}
\begin{itemize}
\raggedright
  \item \emph{Say it:} \emph{suī}
  \item \emph{Say it:} \emph{suī}, once more
  \item \emph{Say it:} say \emph{suī}
  \item \emph{Recall:} say the Marwadi for soap, then the Marwadi for a towel, then say \emph{suī} again
\end{itemize}

\begin{tcolorbox}[breakable,colback=teal!4,colframe=teal!35!black,title={Before you move on}]
What does \mw{सुई} mean? (\textquotedblleft{}A needle\textquotedblright{}.) Say it once more.
\end{tcolorbox}
\section[\texorpdfstring{\mw{डोरो} (ḍoro)}{डोरो (ḍoro)}]{\mw{डोरो} (ḍoro) — thread}

\label{lesson:MW-C124-doro}



\begin{quote}
Before the new one: say the Marwadi for a towel, then the Marwadi for a needle.
\end{quote}

\subsection*{You'll want to know: \mw{डोरो}}
\textbf{\mw{डोरो}} — \emph{ḍoro} — \textquotedblleft{}thread\textquotedblright{}.

\begin{cousinweb}
\textbf{\mw{डोरो}} (\emph{ḍoro}) is a tadbhava, a word inherited from Sanskrit and reshaped by sound change. It comes from \emph{davara}, \textquotedblleft{}string\textquotedblright{}. CDIAL gives no Marwari form in this entry (CDIAL 6225); the match is through Hindi \emph{ḍorā} and Gujarati \emph{doro}.
\end{cousinweb}

\begin{grammarlens}[title={What you've built}]
One more word for this chapter.
\end{grammarlens}

\subsection*{Your turn}
\begin{itemize}
\raggedright
  \item \emph{Say it:} \emph{ḍoro}
  \item \emph{Say it:} \emph{ḍoro}, once more
  \item \emph{Say it:} say \emph{ḍoro}
  \item \emph{Recall:} say the Marwadi for a towel, then the Marwadi for a needle, then say \emph{ḍoro} again
\end{itemize}

\begin{tcolorbox}[breakable,colback=teal!4,colframe=teal!35!black,title={Before you move on}]
What does \mw{डोरो} mean? (\textquotedblleft{}Thread\textquotedblright{}.) Say it once more.
\end{tcolorbox}
